\documentclass{article}
\usepackage{amsmath,amssymb,amsthm}
\usepackage{algorithm,algorithmic}
\usepackage{hyperref}
\usepackage{enumitem}
\newtheorem{theorem}{Theorem}
\newtheorem{lemma}{Lemma}
\newtheorem{assumption}{Assumption}
\newcommand{\E}{\mathbb{E}}
\newcommand{\R}{\mathbb{R}}
\newcommand{\norm}[1]{\left\|#1\right\|}
\newcommand{\ip}[2]{\left\langle #1,#2\right\rangle}
\begin{document}

\section*{Randomized Coordinate Descent with Probability‑Weighted Updates (RCD‑PWU)}

\section*{Mathematical Formulation}
Let $f:\R^{d}\rightarrow\R$ be differentiable and possibly non‑convex.  
For each coordinate $i\in\{1,\dots,d\}$ a probability $p_i>0$ is assigned with
\[
\sum_{i=1}^{d}p_i=1 .
\]

At iteration $t=0,1,\dots$ we draw a random index $i_t$ according to the distribution
$\Pr(i_t=i)=p_i$.  Define the standard basis vector $e_i\in\R^{d}$ with $(e_i)_j=\mathbf 1\{j=i\}$.
The update rule is
\begin{equation}
\boxed{%
w_{t+1}=w_{t}-\frac{\eta}{p_{i_t}}\,e_{i_t}\,\partial_{i_t}f(w_t)
}\label{eq:update}
\end{equation}
where $\eta>0$ is a stepsize (learning rate) and $\partial_i f(w)=\frac{\partial f}{\partial w_i}(w)$.
Because
\[
\E\!\left[\frac{1}{p_{i_t}}e_{i_t}\partial_{i_t}f(w_t)\,\Big|\,w_t\right]
   =\sum_{i=1}^{d}p_i\frac{1}{p_i}e_i\partial_i f(w_t)=\nabla f(w_t),
\]
the stochastic direction in \eqref{eq:update} is an unbiased estimator of the full gradient.

\section*{Pseudocode}
\begin{algorithm}[H]
\caption{RCD‑PWU}
\begin{algorithmic}[1]
\REQUIRE Initial point $w_0\in\R^{d}$, stepsize $\eta>0$, probability vector $p\in\Delta^{d}$.
\FOR{$t=0,1,\dots,T-1$}
    \STATE Sample $i_t\sim p$ (i.e. $\Pr(i_t=i)=p_i$).
    \STATE Compute the partial derivative $g_{i_t}^{(t)}:=\partial_{i_t}f(w_t)$.
    \STATE Update the selected coordinate
    \[
        w_{t+1}=w_{t}-\frac{\eta}{p_{i_t}}\,e_{i_t}\,g_{i_t}^{(t)} .
    \]
\ENDFOR
\RETURN $w_T$ (or any intermediate iterate).
\end{algorithmic}
\end{algorithm}

\section*{Dry Run Example}
Consider the one‑dimensional function $f(w)=(w-3)^2$.
Here $d=1$, $p_1=1$, $\partial_1 f(w)=2(w-3)$.
Choose $w_0=0$ and $\eta=0.1$.

\begin{center}
\begin{tabular}{c|c|c|c}
$t$ & $w_t$ & $\partial_1 f(w_t)$ & $f(w_t)$\\\hline
0 & $0$ & $2(0-3)=-6$ & $9$\\
1 & $0-\frac{0.1}{1}(-6)=0.6$ & $2(0.6-3)=-4.8$ & $(0.6-3)^2=5.76$\\
2 & $0.6-\frac{0.1}{1}(-4.8)=1.08$ & $2(1.08-3)=-3.84$ & $(1.08-3)^2=3.6864$\\
3 & $1.08-\frac{0.1}{1}(-3.84)=1.464$ & $2(1.464-3)=-3.072$ & $(1.464-3)^2=2.3665$\\
\end{tabular}
\end{center}
The objective value decreases at each iteration, illustrating the algorithm.

\newpage
\section*{Assumptions}
\begin{assumption}[Coordinate‑wise $L_i$–smoothness]\label{as:smooth}
For each $i\in\{1,\dots,d\}$ there exists $L_i>0$ such that for all $x\in\R^{d}$ and all $\alpha\in\R$,
\[
\bigl|\partial_i f(x+\alpha e_i)-\partial_i f(x)\bigr|\le L_i |\alpha|.
\]
Equivalently,
\[
f(x+\alpha e_i)\le f(x)+\alpha\,\partial_i f(x)+\frac{L_i}{2}\alpha^{2}.
\]
\end{assumption}

\begin{assumption}[Bounded below]\label{as:lower}
The function $f$ is bounded from below: there exists $f_{\inf}\in\R$ with
$f(w)\ge f_{\inf}$ for all $w\in\R^{d}$.
\end{assumption}

\begin{assumption}[Stepsize]\label{as:stepsize}
The stepsize $\eta$ satisfies
\[
0<\eta\le \min_{i\in[d]}\frac{p_i}{L_i}.
\]
\end{assumption}

\section*{Main Convergence Theorem}
\begin{theorem}\label{thm:main}
Suppose Assumptions \ref{as:smooth}--\ref{as:stepsize} hold.
Define $L_{\max}:=\max_{i}L_i$ and $p_{\min}:=\min_{i}p_i$.
Let $\{w_t\}_{t\ge0}$ be generated by \eqref{eq:update}. Then for any integer $T\ge1$,
\[
\frac{1}{T}\sum_{t=0}^{T-1}\E\!\bigl[\|\nabla f(w_t)\|^{2}\bigr]
   \;\le\;
   \frac{2\bigl(f(w_0)-f_{\inf}\bigr)}{\eta\,T}\,
   \frac{L_{\max}}{p_{\min}} .
\]
Consequently there exists at least one iterate $t\le T-1$ for which
\[
\E\!\bigl[\|\nabla f(w_t)\|^{2}\bigr]
   \;\le\;
   \frac{2\bigl(f(w_0)-f_{\inf}\bigr)}{\eta\,T}\,
   \frac{L_{\max}}{p_{\min}} .
\]
Thus the expected squared norm of the gradient decays at the rate $\mathcal O(1/T)$.
\end{theorem}

\section*{Theoretical Properties}
\begin{itemize}
\item \textbf{Stability.} The bound requires $\eta\le p_i/L_i$ for every coordinate; this guarantees that the quadratic term in the smoothness inequality does not dominate the linear descent term.
\item \textbf{Hyperparameter Sensitivity.} Larger probabilities $p_i$ allow a larger admissible stepsize on coordinate $i$, reflecting the intuition that frequently updated coordinates can move more aggressively.
\item \textbf{Comparison to Uniform CD.} If $p_i=1/d$ for all $i$ and $L_i\equiv L$, the bound reduces to
\[
\frac{2dL\bigl(f(w_0)-f_{\inf}\bigr)}{\eta T},
\]
the classic $\mathcal O(d/L\cdot 1/T)$ rate for uniform random coordinate descent.
\item \textbf{Non‑convexity.} No convexity assumption is needed; the result only relies on smoothness and boundedness below.
\end{itemize}

\section*{Discussion}
The theorem shows that, even in the non‑convex regime, the probability‑weighted update yields a guaranteed decrease in the expected gradient norm. By scaling the step on each coordinate with $1/p_i$, the stochastic direction becomes unbiased, which is the key algebraic device enabling the $\mathcal O(1/T)$ decay without any variance terms.

\newpage
\section*{Preliminaries (Definitions and Lemmas)}
\begin{lemma}[Smoothness inequality for a single coordinate]\label{lem:smooth}
Under Assumption \ref{as:smooth}, for any $x\in\R^{d}$ and any scalar $\alpha$,
\[
f(x+\alpha e_i)\le f(x)+\alpha\,\partial_i f(x)+\frac{L_i}{2}\alpha^{2}.
\]
\end{lemma}
\begin{proof}
The inequality is precisely the integral form of the $L_i$‑Lipschitz condition on the $i$‑th partial derivative. ∎
\end{proof}

\section*{Main Proof (Step‑by‑Step Derivation)}
We prove Theorem \ref{thm:main}.  All expectations are taken with respect to the random index $i_t$ conditioned on the history up to time $t$.

\bigskip
\noindent\textbf{[Step 1]} (Apply coordinate smoothness).  
From Lemma \ref{lem:smooth} with $\alpha=-\frac{\eta}{p_{i_t}}\partial_{i_t}f(w_t)$,
\[
f(w_{t+1})
   \le f(w_t)-\frac{\eta}{p_{i_t}}\bigl(\partial_{i_t}f(w_t)\bigr)^{2}
       +\frac{L_{i_t}}{2}\Bigl(\frac{\eta}{p_{i_t}}\partial_{i_t}f(w_t)\Bigr)^{2}.
\tag{1}
\]

\noindent\textbf{[Step 2]} (Rearrange the quadratic term).  
The last term in (1) can be written as
\[
\frac{L_{i_t}}{2}\frac{\eta^{2}}{p_{i_t}^{2}}\bigl(\partial_{i_t}f(w_t)\bigr)^{2}
   =\frac{\eta}{p_{i_t}}\bigl(\partial_{i_t}f(w_t)\bigr)^{2}
     \underbrace{\frac{L_{i_t}\eta}{2p_{i_t}}}_{\displaystyle\le\frac12}
\]
where the inequality $\frac{L_{i_t}\eta}{2p_{i_t}}\le\frac12$ follows from Assumption \ref{as:stepsize}.

\noindent\textbf{[Step 3]} (One‑step expected descent).  
Insert the bound from Step 2 into (1):
\[
f(w_{t+1})
   \le f(w_t)-\frac{\eta}{p_{i_t}}\bigl(\partial_{i_t}f(w_t)\bigr)^{2}
        +\frac{1}{2}\frac{\eta}{p_{i_t}}\bigl(\partial_{i_t}f(w_t)\bigr)^{2}
   = f(w_t)-\frac{\eta}{2p_{i_t}}\bigl(\partial_{i_t}f(w_t)\bigr)^{2}.
\tag{2}
\]

\noindent\textbf{[Step 4]} (Take conditional expectation).  
Conditioning on $w_t$ and using $\Pr(i_t=i)=p_i$,
\[
\E\!\bigl[f(w_{t+1})\mid w_t\bigr]
   \le f(w_t)-\frac{\eta}{2}\sum_{i=1}^{d}
        \frac{p_i}{p_i}\bigl(\partial_i f(w_t)\bigr)^{2}
   = f(w_t)-\frac{\eta}{2}\,\|\nabla f(w_t)\|^{2}.
\tag{3}
\]

\noindent\textbf{[Step 5]} (Unconditional expectation).  
Taking total expectation of (3) yields
\[
\E\!\bigl[f(w_{t+1})\bigr]
   \le \E\!\bigl[f(w_t)\bigr]-\frac{\eta}{2}\,\E\!\bigl[\|\nabla f(w_t)\|^{2}\bigr].
\tag{4}
\]

\noindent\textbf{[Step 6]} (Telescoping sum).  
Sum (4) from $t=0$ to $T-1$:
\[
\E\!\bigl[f(w_T)\bigr]
   \le f(w_0)-\frac{\eta}{2}\sum_{t=0}^{T-1}
        \E\!\bigl[\|\nabla f(w_t)\|^{2}\bigr].
\tag{5}
\]

\noindent\textbf{[Step 7]} (Use lower bound).  
By Assumption \ref{as:lower}, $f(w_T)\ge f_{\inf}$, therefore
\[
f(w_0)-f_{\inf}
   \ge \frac{\eta}{2}\sum_{t=0}^{T-1}
        \E\!\bigl[\|\nabla f(w_t)\|^{2}\bigr].
\tag{6}
\]

\noindent\textbf{[Step 8]} (Introduce $L_{\max}$ and $p_{\min}$).  
From the definition of $L_{\max}$ and $p_{\min}$ we have for every $i$
\[
\frac{1}{p_i}\le\frac{1}{p_{\min}},\qquad
L_i\le L_{\max}.
\]
Thus $\frac{1}{p_i}\le\frac{L_{\max}}{p_{\min}L_i}$ and consequently
\[
\|\nabla f(w_t)\|^{2}
   =\sum_{i=1}^{d}(\partial_i f(w_t))^{2}
   \le \frac{L_{\max}}{p_{\min}}
        \sum_{i=1}^{d}\frac{p_i}{L_i}(\partial_i f(w_t))^{2}
   \le \frac{L_{\max}}{p_{\min}}
        \sum_{i=1}^{d}\frac{1}{p_i}(\partial_i f(w_t))^{2}.
\]
Applying this inequality inside the sum of (6) gives
\[
\frac{L_{\max}}{p_{\min}}
   \sum_{t=0}^{T-1}
        \E\!\bigl[\|\nabla f(w_t)\|^{2}\bigr]
   \le \frac{2\bigl(f(w_0)-f_{\inf}\bigr)}{\eta}.
\tag{7}
\]

\noindent\textbf{[Step 9]} (Divide by $T$).  
Dividing (7) by $T$ yields the bound stated in Theorem \ref{thm:main}:
\[
\frac{1}{T}\sum_{t=0}^{T-1}\E\!\bigl[\|\nabla f(w_t)\|^{2}\bigr]
   \le \frac{2\bigl(f(w_0)-f_{\inf}\bigr)}{\eta T}\,
      \frac{L_{\max}}{p_{\min}} .
\qquad\blacksquare
\]

\section*{Rate Derivation (With Constants)}
The constant in the rate is 
\[
C:=\frac{2\bigl(f(w_0)-f_{\inf}\bigr)L_{\max}}{\eta p_{\min}} .
\]
Hence the expected squared gradient norm satisfies
\[
\min_{0\le t\le T-1}\E\!\bigl[\|\nabla f(w_t)\|^{2}\bigr]
   \le \frac{C}{T},
\]
which is the optimal $\mathcal O(1/T)$ decay for first‑order methods under only smoothness.

\section*{Special Cases}
\begin{enumerate}[label=(\alph*)]
\item \textbf{Uniform sampling.} If $p_i=1/d$ and $L_i\equiv L$, then $p_{\min}=1/d$,
$L_{\max}=L$, and the bound becomes
\[
\frac{2dL\bigl(f(w_0)-f_{\inf}\bigr)}{\eta T},
\]
exactly the classic rate for uniform random coordinate descent.
\item \textbf{Deterministic cyclic CD.} Setting $p_i=1$ for a single selected coordinate
and $\eta\le 1/L_i$ recovers the deterministic guarantee
$f(w_{t+1})\le f(w_t)-\frac{\eta}{2}(\partial_i f(w_t))^{2}$.
\item \textbf{Heavy‑tailed probabilities.} When some $p_i$ are very small, the admissible
stepsize on those coordinates is proportionally reduced, preventing instability.
\end{enumerate}

\end{document}